\section{完整数据表}

本附录给出正文因篇幅而未完整列出的数据：达标时刻的完整径向剖面、收缩过程中各输出时刻的半径、全部十四个灵敏度算例的逐条结果，以及离散参数收敛研究的原始数值。

\subsection{达标时刻的完整径向剖面}

正文表~\ref{tab:q3C} 与表~\ref{tab:q4C} 按题目要求只给出五个径向位置，表~\ref{tab:profile-full} 给出停时那一刻全部 $21$ 个输出位置的含水率。

\begin{longtable}{ccc}
\caption{停时刻的完整径向含水率剖面（单位：kg/kg）}\label{tab:profile-full}\\
\toprule
到中心距离/cm & 问题三（$t_{\rm dry}=57.5316$~h） & 问题四（$t_{\rm dry}=51.1384$~h） \\
\midrule
\endfirsthead
\multicolumn{3}{l}{\footnotesize 续表~\thetable}\\
\toprule
到中心距离/cm & 问题三（$t_{\rm dry}=57.5316$~h） & 问题四（$t_{\rm dry}=51.1384$~h） \\
\midrule
\endhead
\midrule
\multicolumn{3}{r}{\footnotesize 续下页}\\
\endfoot
\bottomrule
\endlastfoot
0.0 & 0.1500 & 0.1500 \\
0.1 & 0.1499 & 0.1496 \\
0.2 & 0.1496 & 0.1486 \\
0.3 & 0.1492 & 0.1469 \\
0.4 & 0.1485 & 0.1444 \\
0.5 & 0.1477 & 0.1412 \\
0.6 & 0.1466 & 0.1372 \\
0.7 & 0.1454 & 0.1321 \\
0.8 & 0.1439 & 0.1258 \\
0.9 & 0.1422 & 0.1181 \\
1.0 & 0.1402 & 0.1079 \\
1.1 & 0.1379 & 0.0930 \\
1.2 & 0.1353 & 0.0526 \\
1.3 & 0.1324 & --- \\
1.4 & 0.1289 & --- \\
1.5 & 0.1249 & --- \\
1.6 & 0.1202 & --- \\
1.7 & 0.1144 & --- \\
1.8 & 0.1068 & --- \\
1.9 & 0.0956 & --- \\
2.0 & 0.0526 & --- \\
\end{longtable}

问题四一列自 $1.3$~cm 起为破折号，因为该时刻半径已收缩到 $1.2000$~cm，其后的径向位置落在药材之外；$1.2$~cm 处的 $0.0526$ 即当时的表面值。两列对照可见收缩把同样的含水率落差压进了六成的距离内，因此问题四的剖面明显更陡。

\subsection{各输出时刻的药材半径}

问题四的移动边界由附件 2 的实测序列经保形插值给出，表~\ref{tab:radius} 列出正文表~\ref{tab:q4C} 各行对应时刻的半径，供核对该行哪些径向位置仍在药材之内。

\begin{table}[H]
\centering
\caption{问题四各输出时刻的药材半径}\label{tab:radius}
\begin{tabular}{lccccccccc}
\toprule
时间/h & 6 & 12 & 18 & 24 & 30 & 36 & 42 & 48 & 51.1384 \\
\midrule
$R$/cm & 1.374 & 1.248 & 1.214 & 1.204 & 1.201 & 1.200 & 1.200 & 1.200 & 1.200 \\
$R/R_0$ & 0.687 & 0.624 & 0.607 & 0.602 & 0.601 & 0.600 & 0.600 & 0.600 & 0.600 \\
\bottomrule
\end{tabular}
\end{table}

\subsection{灵敏度算例的逐条结果}

正文表~\ref{tab:sens} 按物理量把双侧扰动合并成行，表~\ref{tab:sens-full} 按算例逐条列出，并给出各算例的守恒残差。

\begin{longtable}{llccc}
\caption{全部十四个灵敏度算例（基准为问题三的生产算例）}\label{tab:sens-full}\\
\toprule
算例 & 扰动对象 & $t_{\rm dry}$/h & 相对基准 & $\varepsilon_{\rm mass}$ \\
\midrule
\endfirsthead
\multicolumn{5}{l}{\footnotesize 续表~\thetable}\\
\toprule
算例 & 扰动对象 & $t_{\rm dry}$/h & 相对基准 & $\varepsilon_{\rm mass}$ \\
\midrule
\endhead
\midrule
\multicolumn{5}{r}{\footnotesize 续下页}\\
\endfoot
\bottomrule
\endlastfoot
基准                       & ---            & 57.5316 & ---        & $5.47\times10^{-13}$ \\
$D$ 上调 $10\%$            & 扩散系数       & 52.9223 & $-8.0117\%$  & $5.66\times10^{-13}$ \\
$D$ 下调 $10\%$            & 扩散系数       & 63.1919 & $+9.8387\%$  & $5.38\times10^{-13}$ \\
$h_m$ 上调 $10\%$          & 表面传质系数   & 56.9545 & $-1.0032\%$  & $2.79\times10^{-13}$ \\
$h_m$ 下调 $10\%$          & 表面传质系数   & 58.2672 & $+1.2786\%$  & $2.39\times10^{-13}$ \\
$h$ 上调 $10\%$            & 表面传热系数   & 57.5245 & $-0.0124\%$  & $5.52\times10^{-13}$ \\
$h$ 下调 $10\%$            & 表面传热系数   & 57.5404 & $+0.0152\%$  & $5.63\times10^{-13}$ \\
$C_a$ 上调 $10\%$          & 空气水分浓度   & 57.6522 & $+0.2097\%$  & $5.44\times10^{-13}$ \\
$C_a$ 下调 $10\%$          & 空气水分浓度   & 57.4299 & $-0.1768\%$  & $5.60\times10^{-13}$ \\
$C_{\rm th}=0.165$         & 判据阈值       & 46.0142 & $-20.0193\%$ & $5.48\times10^{-13}$ \\
$C_{\rm th}=0.135$         & 判据阈值       & 76.1486 & $+32.3596\%$ & $5.48\times10^{-13}$ \\
驱动力按体积浓度作差       & 边界写法       & 56.8367 & $-1.2078\%$  & $5.74\times10^{-13}$ \\
温度冻结为空气温度         & 场间耦合       & 57.4462 & $-0.1484\%$  & $5.74\times10^{-13}$ \\
计入潜热汇 $2.4\times10^6$ & 热方程         & 62.9380 & $+9.3972\%$  & $5.20\times10^{-13}$ \\
\end{longtable}

末行的潜热取值为该温区的量级参考值，非题给参数，只用于界定假设（3）的影响，不进入任何交付结果。

\subsection{离散参数收敛研究的原始数值}

表~\ref{tab:conv} 给出网格与时间步收敛研究的逐档结果，正文所引的相对变化均由此表两两相邻计算。

\begin{table}[H]
\centering
\caption{网格与时间步收敛研究的逐档结果}\label{tab:conv}
\begin{tabular}{llcccc}
\toprule
研究项 & 指标 & 第一档 & 第二档 & 第三档 & 第四档 \\
\midrule
网格（问题一） & $N$                  & 20      & 40      & 80      & 160 \\
                & 表面 $C$/(kg/kg)     & 1.51239 & 1.51077 & 1.51037 & 1.51028 \\
网格（问题三） & $N$                  & 20      & 40      & 80      & --- \\
                & $t_{\rm dry}$/h      & 57.5616 & 57.5316 & 57.5237 & --- \\
网格（问题四） & $N$                  & 20      & 40      & 80      & --- \\
                & $t_{\rm dry}$/h      & 51.1664 & 51.1384 & 51.1312 & --- \\
时间步（问题三）& $\Delta t$/s         & 120     & 60      & 30      & 10 \\
                & $t_{\rm dry}$/h      & 57.5783 & 57.5316 & 57.5082 & 57.4926 \\
时间步（问题一）& $\Delta t$/s         & 2       & 1       & 0.5     & 0.25 \\
                & 中心 $T$/°C          & 33.5776 & 33.5764 & 33.5758 & 33.5755 \\
解析解对照     & $\Delta t$/s         & 1       & 0.5     & 0.1     & 0.05 \\
                & 中心 $T$ 偏差/°C     & 3.94e-3 & 1.98e-3 & 4.16e-4 & 2.20e-4 \\
\bottomrule
\end{tabular}

\vspace{0.6em}
\begin{minipage}{0.94\textwidth}
\footnotesize 生产取值为网格 $N=40$、问题三恒温平台段 $\Delta t=60$~s、问题一 $\Delta t=0.01$~s。
时间步（问题一）一行的网格固定为 $N=20$，解析解对照一行的网格固定为 $N=80$、空气温度固定为
平台值 $49.9989$~°C，故这两行的数值与问题一生产算例不属同一算例，不可直接比较。
时间步（问题一）一行按 $2\to1\to0.5\to0.25$~s 逐次减半，式~\eqref{eq:order} 的观测阶 $1.0005$
即由该行后三档算出；此外还加算了减半序列之外的 $\Delta t=0.1$~s 一档，中心 $T$ 为
$33.5753$~°C，与 $0.25$~s 一档相差 $1.8\times10^{-4}$~°C。
\end{minipage}
\end{table}
